\input{preamble.tex}

\title{Лабораторная работа № 16. \\ Базовая защита от атак типа "brute force"}
\author{Сущенко Алина \\ НПИбд-01-23}
\institute{Российский университет дружбы народов имени Патриса Лумумбы}
\date{2025}

\begin{document}

\frame{\titlepage}

\begin{frame}
\frametitle{Цель работы}
\begin{itemize}
    \item Получить навыки работы с программным средством Fail2ban для обеспечения базовой защиты от атак типа "brute force".
\end{itemize}
\end{frame}

\begin{frame}[fragile]
\frametitle{Защита с помощью Fail2ban}
На сервере установили {\tt fail2ban}:
\begin{minted}{bash}
    dnf -y install fail2ban
\end{minted}
    \centering
    \includegraphics[width=\textwidth]{../images/01.png}
    \captionof{figure}{Запуск сервиса \texttt{fail2ban}.}
\end{frame}

\begin{frame}
\frametitle{Защита с помощью Fail2ban}
    \centering
    \includegraphics[width=\textwidth]{../images/02.png}
    \captionof{figure}{Просмотр журнала событий \texttt{fail2ban}.}
\end{frame}

\begin{frame}[fragile]
\frametitle{Защита с помощью Fail2ban}
\begin{itemize}
\item Создали файл с новой конфигурацией {\tt fail2ban}:
\begin{minted}{bash}
    touch /etc/fail2ban/jail.d/customisation.local
\end{minted}
\item Задали время блокирования на 1 час (время задаётся в секундах):
\begin{minted}{bash}
      [DEFAULT]
      bantime = 3600
\end{minted}
  \item включили защиту SSH:
\begin{minted}{bash}
      [sshd]
      port = ssh,2022
      enabled = true
      [sshd-dos]
      filter = sshd
      enabled = true
      [selinux-ssh]
      enabled = true
\end{minted}
\end{itemize}
\end{frame}

\begin{frame}
\frametitle{Защита с помощью Fail2ban}
    \centering
    \includegraphics[width=\textwidth]{../images/03.png}
    \captionof{figure}{Просмотр журнала событий после первоначальной настройки \texttt{fail2ban}.}
\end{frame}

\begin{frame}[fragile]
\frametitle{Защита с помощью Fail2ban}
В файле {\tt /etc/fail2ban/jail.d/customisation.local} включите защиту HTTP:
\begin{minted}[fontsize=\footnotesize]{bash}
    # HTTP servers
    [apache-auth]
    enabled = true
    [apache-badbots]
    enabled = true
    [apache-noscript]
    enabled = true
    [apache-overflows]
    enabled = true
    [apache-nohome]
    enabled = true
    [apache-botsearch]
    enabled = true
    [apache-fakegooglebot]
    enabled = true
    [apache-modsecurity]
    enabled = true
    [apache-shellshock]
    enabled = true
\end{minted}
\end{frame}

\begin{frame}
\frametitle{Защита с помощью Fail2ban}
    \centering
    \includegraphics[width=\textwidth]{../images/04.png}
    \captionof{figure}{Логи, связанные с защитой HTTP (Часть 1).}
\end{frame}

\begin{frame}
\frametitle{Защита с помощью Fail2ban}
    \centering
    \includegraphics[width=\textwidth]{../images/05.png}
    \captionof{figure}{Логи, связанные с защитой HTTP (Часть 2).}
\end{frame}

\begin{frame}[fragile]
\frametitle{Защита с помощью Fail2ban}
В файле {\tt /etc/fail2ban/jail.d/customisation.local} включите защиту почты:
\begin{minted}{bash}
    #
    # Mail servers
    #
    [postfix]
    enabled = true
    [postfix-rbl]
    enabled = true
    [dovecot]
    enabled = true
    [postfix-sasl]
    enabled = true
\end{minted}
\end{frame}

\begin{frame}
\frametitle{Защита с помощью Fail2ban}
  \centering
  \includegraphics[width=\textwidth]{../images/06.png}
  \captionof{figure}{Логи, связанные с защитой почты.}
\end{frame}


\begin{frame}
\frametitle{Проверка работы Fail2ban}
    \centering
    \includegraphics[width=\textwidth]{../images/07.png}
    \captionof{figure}{Просмотр статуса \texttt{fail2ban}.}
\end{frame}

\begin{frame}
\frametitle{Проверка работы Fail2ban}
    \centering
    \includegraphics[width=\textwidth]{../images/08.png}
    \captionof{figure}{Просмотр статуса защиты SSH.}
\end{frame}

\begin{frame}
\frametitle{Проверка работы Fail2ban}
    \centering
    \includegraphics[width=\textwidth]{../images/09.png}
    \captionof{figure}{Установка максимального количества ошибок для SSH.}
\end{frame}

\begin{frame}
\frametitle{Проверка работы Fail2ban}
    \centering
    \includegraphics[width=\textwidth]{../images/10.png}
    \captionof{figure}{Попытка зайти на сервер через SSH с неправильным паролем.}
\end{frame}

\begin{frame}
\frametitle{Проверка работы Fail2ban}
    \centering
    \includegraphics[width=\textwidth]{../images/11.png}
    \captionof{figure}{Просмотр статуса защиты SSH.}
\end{frame}

\begin{frame}
\frametitle{Проверка работы Fail2ban}
    \centering
    \includegraphics[width=\textwidth]{../images/12.png}
    \captionof{figure}{Разблокировка ip-адреса клиента.}
\end{frame}

\begin{frame}
\frametitle{Проверка работы Fail2ban}
    \centering
    \includegraphics[width=\textwidth]{../images/13.png}
    \captionof{figure}{Проверка снятия блокировки.}
\end{frame}

\begin{frame}[fragile]
\frametitle{Проверка работы Fail2ban}
На сервере внесли изменение в конфигурационный файл {\tt /etc/fail2ban/jail.d/customisation.local}, добавив в раздел по умолчанию игнорирование адреса клиента
\begin{minted}{bash}
    [DEFAULT]
    bantime = 3600
    ingoreip = 127.0.0.1/8 192.168.1.30
\end{minted}

    \centering
    \includegraphics[width=\textwidth]{../images/14.png}
    \captionof{figure}{Логи журнала об игнорировании адреса клиента.}
\end{frame}

\begin{frame}
\frametitle{Проверка работы Fail2ban}
    \centering
    \includegraphics[width=\textwidth]{../images/15.png}
    \captionof{figure}{Демонстрация игнорирования безуспешных попыток подключения по SSH.}
\end{frame}

\begin{frame}
\frametitle{Проверка работы Fail2ban}
    \centering
    \includegraphics[width=\textwidth]{../images/16.png}
    \captionof{figure}{Просмотр информации о защите SSH.}
\end{frame}

\begin{frame}[fragile]
\frametitle{Внесение изменений в настройки внутреннего окружения виртуальных машин}
\begin{minted}{bash}
    #!/bin/bash
    echo "Provisioning script $0"
    echo " Install needed packages"
    dnf -y install fail2ban
    echo "Copy configuration files"
    cp -R /vagrant/provision/server/protect/etc/* /etc
    restorecon -vR /etc
    echo "Start fail2ban service"
    systemctl enable fail2ban
    systemctl start fail2ban
\end{minted}
\end{frame}

\begin{frame}[fragile]
\frametitle{Внесение изменений в настройки внутреннего окружения виртуальных машин}
\begin{minted}{bash}
    server.vm.provision "server protect",
        type: "shell",
        preserve_order: true,
        path: "provision/server/protect.sh"
\end{minted}
\end{frame}

\begin{frame}
\frametitle{Контрольные вопросы}
1. Поясните принцип работы Fail2ban.
\begin{itemize}
    \item Fail2ban сканирует логи служб (SSH, Apache, Postfix и др.) на предмет неудачных попыток аутентификации.
    \item При превышении заданного количества попыток (\texttt{maxretry}) в течение определённого времени (\texttt{findtime}) система добавляет IP-адрес нарушителя в черный список.
    \item Блокировка осуществляется через встроенные механизмы (iptables, firewalld) на заданное время (\texttt{bantime}).
    \item После истечения времени блокировки IP-адрес автоматически разблокируется.
\end{itemize}

2. Настройки какого файла более приоритетны: jail.conf или jail.local?
\begin{itemize}
    \item Настройки файла \texttt{jail.local} имеют более высокий приоритет.
    \item Fail2ban читает сначала \texttt{jail.conf}, затем \texttt{jail.local}, перезаписывая совпадающие параметры.
    \item Рекомендуется редактировать \texttt{jail.local} для внесения изменений, оставляя \texttt{jail.conf} в исходном состоянии.
\end{itemize}
\end{frame}

\begin{frame}
\frametitle{Контрольные вопросы (продолжение)}
3. Как настроить оповещение администратора при срабатывании Fail2ban?
\begin{itemize}
    \item Использовать действие \texttt{action\_mwl} в конфигурации jail: \texttt{action = \%\texttt{(action\_mwl)}s}
    \item Настроить параметры почты в секции \texttt{[DEFAULT]}:
    \begin{itemize}
        \item \texttt{mta = sendmail} (или postfix, exim)
        \item \texttt{destemail = admin@example.com}
        \item \texttt{sender = fail2ban@example.com}
    \end{itemize}
    \item Для более сложных уведомлений можно создать собственное действие в \texttt{/etc/fail2ban/action.d/}
\end{itemize}
\end{frame}
4. Поясните построчно настройки по умолчанию в конфигурационном файле /etc/fail2ban/jail.conf, относящиеся к веб-службе.
\begin{itemize}
\begin{frame}
    \item \texttt{[apache-auth]} - защита от перебора паролей аутентификации Apache
    \item \texttt{enabled = false} - по умолчанию отключено
    \item \texttt{port = http,https} - порты для мониторинга
    \item \texttt{filter = apache-auth} - используемый фильтр
    \item \texttt{logpath = /var/log/apache*/*error*.log} - путь к логам
    \item \texttt{maxretry = 6} - максимальное количество попыток
    \item \texttt{bantime = 600} - время блокировки (600 секунд)
\end{itemize}
\end{frame}

\begin{frame}
\frametitle{Контрольные вопросы (продолжение)}
5. Поясните построчно настройки по умолчанию в конфигурационном файле /etc/fail2ban/jail.conf, относящиеся к почтовой службе.
\begin{itemize}
    \item \texttt{[postfix]} - защита сервера Postfix
    \item \texttt{enabled = false} - по умолчанию отключено
    \item \texttt{port = smtp,ssmtp,submission} - порты SMTP
    \item \texttt{filter = postfix} - фильтр для анализа логов Postfix
    \item \texttt{logpath = /var/log/mail.log} - путь к почтовым логам
    \item \texttt{maxretry = 5} - максимальное количество попыток
    \item \texttt{[dovecot]} - аналогично для Dovecot IMAP/POP3 сервера
    \item \texttt{port = pop3,pop3s,imap,imaps} - порты IMAP/POP3
    \item \texttt{filter = dovecot} - фильтр для логов Dovecot
\end{itemize}
\end{frame}

\begin{frame}
6. Какие действия может выполнять Fail2ban при обнаружении атакующего IP-адреса? Где можно посмотреть описание действий для последующего использования в настройках Fail2ban?
\begin{itemize}
    \item Основные действия:
    \begin{itemize}
        \item Блокировка через iptables/firewalld
        \item Отправка email уведомлений
        \item Выполнение произвольных команд
        \item Отправка сообщений в syslog
        \item Уведомление через Telegram, Slack и др.
    \end{itemize}
    \item Описание действий находится в директории \texttt{/etc/fail2ban/action.d/}
    \item Например, \texttt{iptables.conf}, \texttt{mail.conf}, \texttt{sendmail.conf}
\end{itemize}
\end{frame}

\begin{frame}
\frametitle{Контрольные вопросы (продолжение)}
7. Как получить список действующих правил Fail2ban?
\begin{itemize}
    \item Команда: \texttt{fail2ban-client status}
    \item Для конкретной службы: \texttt{fail2ban-client status sshd}
    \item Просмотр через iptables: \texttt{iptables -L -n} или \texttt{iptables -S}
    \item Для firewalld: \texttt{firewall-cmd --list-all}
\end{itemize}

8. Как получить статистику заблокированных Fail2ban адресов?
\begin{itemize}
    \item \texttt{fail2ban-client status <jail-name>} - показывает количество заблокированных IP
    \item \texttt{fail2ban-client get <jail-name> banned} - список заблокированных IP
    \item \texttt{journalctl -u fail2ban} - журнал событий Fail2ban
    \item \texttt{/var/log/fail2ban.log} - основной лог-файл
\end{itemize}
\end{frame}

\begin{frame}
9. Как разблокировать IP-адрес?
\begin{itemize}
    \item Команда: \texttt{fail2ban-client set <jail> unbanip <IP>}
    \item Пример: \texttt{fail2ban-client set sshd unbanip 192.168.1.100}
    \item Удаление правила iptables вручную: \texttt{iptables -D fail2ban-sshd -s <IP> -j DROP}
    \item Перезагрузка службы: \texttt{systemctl restart fail2ban} (разблокирует все адреса)
\end{itemize}
\end{frame}

\begin{frame}
\frametitle{Выводы}
\begin{itemize}
    \item Во время выполнения лабораторной работы получили навыки настройки обеспечения базовой защиты от атак типа "brute force" с помощью программы Fail2ban.
\end{itemize}
\end{frame}
\end{document}
